% !TeX root = ../papers/nc-p3.tex
\section{階数14の全幾何剰余点と円分倍率の正則性}
\label{sec:rank-fourteen-geometric-reduction}
\begin{theorem}\label{thm:rank-fourteen-geometric-reduction}
標数零の階数14の窓が離散付値環上の通常表示を持ち、その還元が標数101の通常点ならば、
分数体上で標準Koszul左右AS正則である。剰余座標 $t,s$ の拡大次数に制限はない。
特に、有理既約な原始整数多項式 $P(t)$ の101での還元が通常な幾何根を一つ持てば、
$P$ の全共役の上の全ての窓が正則である。この根には単根性を仮定しない。
\end{theorem}
\begin{proof}
$k=\mathbb F_{101}$、$R=k[t]$、$C=2(t-1)(t+1)(t+3)$、
$S=R[Y]/(Y^2-C)$ と置く。$i=10,91$ のそれぞれに対し、前節の四式を
分母が可逆な点で使う。二次関係と四次中心式は一行ずつ共通分母を払って $S$ 上に記す。
以下の計算は全ての $t$ を一変数のまま行う。

始点 $p$ を固定する。長さ $n$、文字 $h$ の語に $S$ の値を与え、
関係を零にする線型汎関数を考える。$R$ へ係数を制限すると、空語の値の空間は
階数2で、$Y$ の作用は $M_0=\left(\begin{smallmatrix}0&C\\1&0\end{smallmatrix}\right)$
である。最後の文字 $b$ を指定した制限写像を $J_{n,h}^{b}$ と書く。
長さ $n-1$ の四つの空間の直和の元を $(v_b)_b$ とする。
末尾に置く文字 $g$ の二次関係 $\sum c_{ab}x_ax_b$ に対し、条件は
\[
 \sum_{a,b} c_{ab}(M_{n-2,h+g})\,J_{n-1,h+b}^{a}\,v_b=0 .
\]
ここで $c=u+vY$ に対する $c(M)$ は $uI+vM$ である。
長さ4では各中心式にも同様に三つの制限写像の積を用いる。
後続の長さでは二次関係だけを加える。この条件だけでも元の商の双対を含み、
分数体で四式が中心なら前節の交換則により全条件となる。

得られた行列を $A$ とする。$R$ は単項イデアル整域なので、
多項式上の核の基底を並べた $J$ が取れる。その四つの行ブロックが次の制限写像となる。
$Y$ 作用を直和した行列 $M'$ に対し $JM=M'J$ を解き、全成分で確認する。
$A$ の列加群の基底を並べた行列を $B$、その行加群の基底を並べた正方行列を $D$ とする。
$r=\operatorname{rank}A$ とすれば、$A$ の $r$ 次小行列式のイデアルは $(\det D)$ である。
実際、生成列と基底列は互いに $R$ 線型結合で書け、行についても同様なので
Cauchy--Binet公式を両方向に適用できる。したがって $\det D$ が単元の局所環では
核と特殊化が交換する。一変数の被約Gröbner基底は先頭成分が相異なるため、
ここで必要な各部分加群の自由基底を与える。

\path{computations/rank_fourteen_dual.py} は始点 $0,1$ と各 $i$ で
長さ2から13の全四文字についてこの計算を行う。各長さの核の総階数は
命題\ref{prop:fixed-residual-regularity}の商次元の二倍である。
$C$ は $k(t)$ で平方ではなく、一般核は $k(t)(Y)$ 上の空間である。
通常点では二根が相異なる。よって定階数の開集合で $Y$ の二つの固有空間は
それぞれ一般階数の半分を持つ（有限エタール被覆上で各直和因子の階数は局所一定）。
二位置移動で全四始点が扱え、共通の四式の選択には前節と同じ開集合の交叉を使える。

$\det D$ の既約因子と四式および関係の分母を合わせ、その全根も調べる。
\path{computations/verify_rank_fourteen_geometric_reduction.py} は
各通常既約因子の次数を $d$ として $\mathbb F_{101^d}$ に根を取り、
必要なら二次拡大して $Y$ を取る。Frobenius共役と二位置移動により、
各 $i$ で一つの $t,Y$ と全四始点を検査すれば足りる。各点で経路長4,5の次元
35,56、中心の文字次元 $(3,2,0,0)$ を直接求め直し、その中心から選んだ
四式の長さ13までの商次元が所定の列に一致することを確かめる。

分母も行列式も消えない還元では有理式で指定した中心元が分数体上で使え、
還元で得た商次元は分数体での上界を与える。例外因子での直接検査では
定階数の経路・中心を経由して剰余切断を持ち上げる。
いずれも一般正則点からの全次数の下界を合わせれば
命題\ref{prop:fixed-residual-regularity}の議論が適用される。
最後に $P$ を射影直線上で原始化して閉包を取れば、101でのその幾何根は
正規化上の付値に支配される。そこで数体を必要なだけ拡大すれば $i,Y$ の全ての選択ができ、
同じ結論が共役で全ての根へ移る。
\end{proof}

\begin{corollary}\label{cor:rank-fourteen-cyclotomic}
階数14の二成分で、$t$ が1の冪根である全ての通常窓は、標数零で
標準Koszul左右AS正則である。
\end{corollary}
\begin{proof}
$t$ の位数を $101^a m$、$101\nmid m$ とする。101における共役の還元は
全ての原始 $m$ 乗根を尽くす。非通常な有限 $t$ は101で相異なる13点であり、
そのうちに現れる乗法位数は $1,2,25,51,100,3400,10200$ である。
この集合の $m>2$ は全て $\varphi(m)>13$ を満たす。従って $m>2$ なら
原始根の一つが通常であり、定理が適用される。

$m=1,2$ で $a=0$ のとき $t=\pm1$ は非通常である。$a>0$ のときは
素数109を使う。$109$ の101を法とする位数は100なので、原始 $101^a m$ 乗根の
109での剰余次数は100の倍数となる。$i=33$ として上の多項式核と行列式を
二始点で計算すると、分母と非通常因子を含む全例外因子の次数は100未満である。
\path{computations/verify_rank_fourteen_cyclotomic.py} が両素数での因数分解と位数を再現する。
$\mathbb Q(t)$ は奇数導手の円分体で2で不分岐だから $\mathbb Q(i)$ と交わるのは
$\mathbb Q$ だけである。従って $i$ のどちらの選択にもこの還元を使え、
$Y$ の両根は上と同じ核計算で扱われる。
\end{proof}
